% GENERATED FILE. Edit canonical lessons, then run npm run generate:books.
% canonical-source-hash: fnv1a64:9e61edac581a62ce
% canonical-lessons: ML-C373-akramanam, ML-C373-bhisani, ML-C373-akramikkuka, ML-C373-kollayatikkuka, ML-C373-takarkkuka

\chapter{An Attack}
\label{ch:ml-an-attack}
\hlchaptermodality{373}

\begin{chapteropening}
\textbf{By the end of this chapter:} \emph{I can say an attack, a threat, to attack, to loot and to smash.}

Complete the last lesson of chapter 373: I can say an attack, a threat, to attack, to loot and to smash.
\end{chapteropening}

\section[\texorpdfstring{\ml{ആക്രമണം} (ākramaṇaṁ)}{ആക്രമണം (ākramaṇaṁ)}]{\ml{ആക്രമണം} (ākramaṇaṁ) — an attack}

\label{lesson:ML-C373-akramanam}



\begin{quote}
Before the new one: say the Malayalam for to trick, then the Malayalam for to fool.
\end{quote}

\subsection*{You'll want to know: \ml{ആക്രമണം}}
\textbf{\ml{ആക്രമണം}} — \emph{ākramaṇaṁ} — \textquotedblleft{}an attack\textquotedblright{}.

\begin{quote}
\ml{ആനയുടെ} \ml{ആക്രമണത്തിൽ} \ml{ഒരു} \ml{വീട്} \ml{തകർന്നു}.
\end{quote}

\begin{quote}
\emph{ānayuṭe ākramaṇattil oru vīṭŭ takarnnu.}
\end{quote}

\begin{quote}
A house was wrecked in an elephant's attack.
\end{quote}

Also of animals: an elephant or a dog attack.

\begin{cousinweb}
\textbf{\ml{ആക്രമണം}} (\emph{ākramaṇaṁ}) is borrowed from Sanskrit \emph{ākramaṇa}, \textquotedblleft{}stepping on, approaching\textquotedblright{} (CDIAL 1016), from \emph{ā-kram}, \textquotedblleft{}to step upon\textquotedblright{}: an attack steps onto someone's ground.
\end{cousinweb}

\begin{grammarlens}[title={What you've built}]
One more word for telling what happened.
\end{grammarlens}

\subsection*{Your turn}
\begin{itemize}
\raggedright
  \item \emph{Say it:} \emph{ākramaṇaṁ}
  \item \emph{Say it:} \emph{ākramaṇaṁ}, once more
  \item \emph{Say it:} \emph{ānayuṭe ākramaṇattil oru vīṭŭ takarnnu.} — A house was wrecked in an elephant's attack.
  \item \emph{Recall:} say the Malayalam for to trick, then the Malayalam for to fool, then say \emph{ākramaṇaṁ} again
\end{itemize}

\begin{tcolorbox}[breakable,colback=teal!4,colframe=teal!35!black,title={Before you move on}]
What does \ml{ആക്രമണം} mean? (\textquotedblleft{}An attack\textquotedblright{}.) Say it once more.
\end{tcolorbox}
\section[\texorpdfstring{\ml{ഭീഷണി} (bhīṣaṇi)}{ഭീഷണി (bhīṣaṇi)}]{\ml{ഭീഷണി} (bhīṣaṇi) — a threat}

\label{lesson:ML-C373-bhisani}



\begin{quote}
Before the new one: say the Malayalam for to fool, then the Malayalam for an attack.
\end{quote}

\subsection*{You'll want to know: \ml{ഭീഷണി}}
\textbf{\ml{ഭീഷണി}} — \emph{bhīṣaṇi} — \textquotedblleft{}a threat\textquotedblright{}.

\begin{quote}
\ml{അവന്} \ml{ഫോണിൽ} \ml{ഒരു} \ml{ഭീഷണി} \ml{വന്നു}.
\end{quote}

\begin{quote}
\emph{avanŭ phōṇil oru bhīṣaṇi vannu.}
\end{quote}

\begin{quote}
He received a threat over the phone.
\end{quote}

A threat to something takes -\ml{ക്ക്}: \ml{ആരോഗ്യത്തിന്} \ml{ഭീഷണി}, a threat to health.

\begin{cousinweb}
\textbf{\ml{ഭീഷണി}} (\emph{bhīṣaṇi}) comes from Sanskrit \emph{bhīṣaṇa}, \textquotedblleft{}terrifying, frightening\textquotedblright{} (Monier-Williams), a causative form of \emph{bhī}, \textquotedblleft{}to fear\textquotedblright{}.
\end{cousinweb}

\begin{grammarlens}[title={What you've built}]
One more word for telling what happened.
\end{grammarlens}

\subsection*{Your turn}
\begin{itemize}
\raggedright
  \item \emph{Say it:} \emph{bhīṣaṇi}
  \item \emph{Say it:} \emph{bhīṣaṇi}, once more
  \item \emph{Say it:} \emph{avanŭ phōṇil oru bhīṣaṇi vannu.} — He received a threat over the phone.
  \item \emph{Recall:} say the Malayalam for to fool, then the Malayalam for an attack, then say \emph{bhīṣaṇi} again
\end{itemize}

\begin{tcolorbox}[breakable,colback=teal!4,colframe=teal!35!black,title={Before you move on}]
What does \ml{ഭീഷണി} mean? (\textquotedblleft{}A threat\textquotedblright{}.) Say it once more.
\end{tcolorbox}
\section[\texorpdfstring{\ml{ആക്രമിക്കുക} (ākramikkuka)}{ആക്രമിക്കുക (ākramikkuka)}]{\ml{ആക്രമിക്കുക} (ākramikkuka) — to attack}

\label{lesson:ML-C373-akramikkuka}



\begin{quote}
Before the new one: say the Malayalam for an attack, then the Malayalam for a threat.
\end{quote}

\subsection*{You'll want to know: \ml{ആക്രമിക്കുക}}
\textbf{\ml{ആക്രമിക്കുക}} — \emph{ākramikkuka} — \textquotedblleft{}to attack\textquotedblright{}.

\begin{quote}
\ml{ഒരു} \ml{നായ} \ml{കുട്ടിയെ} \ml{ആക്രമിച്ചു}.
\end{quote}

\begin{quote}
\emph{oru nāya kuṭṭiye ākramiccu.}
\end{quote}

\begin{quote}
A dog attacked the child.
\end{quote}

Past: \ml{ആക്രമിച്ചു}, \emph{ākramiccu}.

The one attacked takes -\ml{എ}: \ml{അവനെ} \ml{ആക്രമിച്ചു}.

\begin{cousinweb}
\textbf{\ml{ആക്രമിക്കുക}} (\emph{ākramikkuka}) is built on the Sanskrit verb \emph{ā-kram}, \textquotedblleft{}to step upon\textquotedblright{}, whose noun \emph{ākramaṇa} (CDIAL 1016) gave \ml{ആക്രമണം} (\emph{ākramaṇaṁ}), with the Malayalam verb ending -\ml{ിക്കുക}.
\end{cousinweb}

\begin{grammarlens}[title={What you've built}]
One more word for telling what happened.
\end{grammarlens}

\subsection*{Your turn}
\begin{itemize}
\raggedright
  \item \emph{Say it:} \emph{ākramikkuka}
  \item \emph{Say it:} \emph{ākramikkuka}, once more
  \item \emph{Say it:} \emph{oru nāya kuṭṭiye ākramiccu.} — A dog attacked the child.
  \item \emph{Recall:} say the Malayalam for an attack, then the Malayalam for a threat, then say \emph{ākramikkuka} again
\end{itemize}

\begin{tcolorbox}[breakable,colback=teal!4,colframe=teal!35!black,title={Before you move on}]
What does \ml{ആക്രമിക്കുക} mean? (\textquotedblleft{}To attack\textquotedblright{}.) Say it once more.
\end{tcolorbox}
\section[\texorpdfstring{\ml{കൊള്ളയടിക്കുക} (koḷḷayaṭikkuka)}{കൊള്ളയടിക്കുക (koḷḷayaṭikkuka)}]{\ml{കൊള്ളയടിക്കുക} (koḷḷayaṭikkuka) — to loot; to plunder}

\label{lesson:ML-C373-kollayatikkuka}



\begin{quote}
Before the new one: say the Malayalam for a threat, then the Malayalam for to attack.
\end{quote}

\subsection*{You'll want to know: \ml{കൊള്ളയടിക്കുക}}
\textbf{\ml{കൊള്ളയടിക്കുക}} — \emph{koḷḷayaṭikkuka} — \textquotedblleft{}to loot; to plunder\textquotedblright{}.

\begin{quote}
\ml{കള്ളന്മാർ} \ml{രാത്രി} \ml{കട} \ml{കൊള്ളയടിച്ചു}.
\end{quote}

\begin{quote}
\emph{kaḷḷanmār rātri kaṭa koḷḷayaṭiccu.}
\end{quote}

\begin{quote}
The robbers looted the shop at night.
\end{quote}

Past: \ml{കൊള്ളയടിച്ചു}, \emph{koḷḷayaṭiccu}.

\ml{കൊള്ള}, plunder, + \ml{അടിക്കുക}, to strike.

\begin{cousinweb}
\textbf{\ml{കൊള്ളയടിക്കുക}} (\emph{koḷḷayaṭikkuka}) is native Dravidian: \ml{കൊള്ള}, \textquotedblleft{}plunder\textquotedblright{}, which DEDR 2151 lists with \emph{koḷ}, \textquotedblleft{}to hold, to take\textquotedblright{}, and \ml{അടിക്കുക} (\emph{aṭikkuka}), \textquotedblleft{}to strike\textquotedblright{} (DEDR 77).
\end{cousinweb}

\begin{grammarlens}[title={What you've built}]
One more word for telling what happened.
\end{grammarlens}

\subsection*{Your turn}
\begin{itemize}
\raggedright
  \item \emph{Say it:} \emph{koḷḷayaṭikkuka}
  \item \emph{Say it:} \emph{koḷḷayaṭikkuka}, once more
  \item \emph{Say it:} \emph{kaḷḷanmār rātri kaṭa koḷḷayaṭiccu.} — The robbers looted the shop at night.
  \item \emph{Recall:} say the Malayalam for a threat, then the Malayalam for to attack, then say \emph{koḷḷayaṭikkuka} again
\end{itemize}

\begin{tcolorbox}[breakable,colback=teal!4,colframe=teal!35!black,title={Before you move on}]
What does \ml{കൊള്ളയടിക്കുക} mean? (\textquotedblleft{}To loot; to plunder\textquotedblright{}.) Say it once more.
\end{tcolorbox}
\section[\texorpdfstring{\ml{തകർക്കുക} (takarkkuka)}{തകർക്കുക (takarkkuka)}]{\ml{തകർക്കുക} (takarkkuka) — to smash; to wreck}

\label{lesson:ML-C373-takarkkuka}



\begin{quote}
Before the new one: say the Malayalam for to attack, then the Malayalam for to loot.
\end{quote}

\subsection*{You'll want to know: \ml{തകർക്കുക}}
\textbf{\ml{തകർക്കുക}} — \emph{takarkkuka} — \textquotedblleft{}to smash; to wreck\textquotedblright{}.

\begin{quote}
\ml{കാറ്റ്} \ml{ഗ്രാമത്തിലെ} \ml{വീടുകൾ} \ml{തകർത്തു}.
\end{quote}

\begin{quote}
\emph{kāṟṟŭ grāmattile vīṭukaḷ takarttu.}
\end{quote}

\begin{quote}
The wind wrecked the houses of the village.
\end{quote}

Past: \ml{തകർത്തു}, \emph{takarttu}.

The partner of \ml{തകരുക}, to collapse: one smashes, the other is smashed.

\begin{cousinweb}
\textbf{\ml{തകർക്കുക}} (\emph{takarkkuka}) is native Dravidian: DEDR 2999 lists Malayalam \emph{takarkka}, \textquotedblleft{}to smash, crush, demolish\textquotedblright{}, the active partner of \ml{തകരുക} (\emph{takaruka}).
\end{cousinweb}

\begin{grammarlens}[title={What you've built}]
One more word for telling what happened.
\end{grammarlens}

\subsection*{Your turn}
\begin{itemize}
\raggedright
  \item \emph{Say it:} \emph{takarkkuka}
  \item \emph{Say it:} \emph{takarkkuka}, once more
  \item \emph{Say it:} \emph{kāṟṟŭ grāmattile vīṭukaḷ takarttu.} — The wind wrecked the houses of the village.
  \item \emph{Recall:} say the Malayalam for to attack, then the Malayalam for to loot, then say \emph{takarkkuka} again
\end{itemize}

\begin{tcolorbox}[breakable,colback=teal!4,colframe=teal!35!black,title={Before you move on}]
What does \ml{തകർക്കുക} mean? (\textquotedblleft{}To smash; to wreck\textquotedblright{}.) Say it once more.
\end{tcolorbox}
